\chapter{Masked Diffusion Language Models}
\label{ch:masked-diffusion}

%% Status: WRITE NOW — L&Z reading complete (2026-04-11).
%% Depends on: ch3; L&Z decomposition is the core of sec:lz-decomposition
%%
%% Format: ~1200--1500 words + equations.

%% ---------------------------------------------------------------
\section{MDLM: The SUBS Parameterization}
\label{sec:mdlm-subs}

%% Cover (based on MDLM reading):
%%   - The SUBS (Substitution) parameterization: p_theta(x | z_t) directly
%%   - Training objective: masked language modelling loss
%%   - Sampling algorithm: sequential unmasking from z_T to z_0
%%   - Easy-first unmasking: why high-confidence positions are unmasked first
%%   - When the factorisation approximation hurts most

%% ---------------------------------------------------------------
\section{Remasking Strategies: ReMDM}
\label{sec:remdm}

%% Cover (based on ReMDM reading):
%%   - The remasking backward process: why remasking improves quality
%%   - Four schedule variants: conf, loop, cap, rescale
%%   - Why confidence-based remasking creates OOD states at high T
%%   - Connection between remasking budget and number of NFEs
%%   - Key distinction: remasking defers tokens (predictor modification)
%%                      correctors add extra MCMC steps (our focus)

%% ---------------------------------------------------------------
\section{The L\&Z Error Decomposition}
\label{sec:lz-decomposition}

%% Cover (based on L&Z reading — complete this after finishing L&Z):
%%   - The KL decomposition: \KL{\ptrue}{\palg} \leq \Elearn + \Efact
%%   - \Elearn: controlled by training (neural denoiser quality)
%%   - \Efact: controlled at inference time by the choice of schedule and correctors
%%   - The EB-Sampler: minimises \Efact among unmasking-only strategies
%%   - Section 6 of L&Z: "extending to corrector steps" as future work
%%   - This thesis aims to answer that open question

%% ---------------------------------------------------------------
\section{The Corrector Problem}
\label{sec:corrector-problem}

%% Position the thesis within this framework:
%%   - What a corrector is: extra Markov step after token commitment
%%   - How a corrector appears in the total error bound
%%   - The open question: can an informed corrector reduce \Efact below the EB-Sampler bound?
%%   - Preview of the candidate directions (Gap D, B/C, E)
